\documentclass[a4paper, serif, 10pt]{moderncv}

%% moderncv
% style and color
\moderncvstyle{banking}
\moderncvcolor{black}

% renew moderncv command to adapt for CJK colon
\renewcommand*{\cvitem}[3][.25em]{%
  \ifstrempty{#2}{}{\hintstyle{#2}：}{#3}%
  \par\addvspace{#1}}

\renewcommand*{\cvitemwithcomment}[4][.25em]{%
  \savebox{\cvitemwithcommentmainbox}{\ifstrempty{#2}{}{\hintstyle{#2}：}#3}%
  \setlength{\cvitemwithcommentmainlength}{\widthof{\usebox{\cvitemwithcommentmainbox}}}%
  \setlength{\cvitemwithcommentcommentlength}{\maincolumnwidth-\separatorcolumnwidth-\cvitemwithcommentmainlength}%
  \begin{minipage}[t]{\cvitemwithcommentmainlength}\usebox{\cvitemwithcommentmainbox}\end{minipage}%
  \hfill% fill of \separatorcolumnwidth
  \begin{minipage}[t]{\cvitemwithcommentcommentlength}\raggedleft\small\itshape#4\end{minipage}%
  \par\addvspace{#1}}

%% page layout/margins
\usepackage[top=2.5cm, bottom=2.5cm, left=1.5cm, right=1.5cm]{geometry}

\nopagenumbers{}

%% Babel config for English language
\usepackage[english]{babel}

%% fontspec
\usepackage{fontspec}

\IfFontExistsTF{Linux Libertine}{
  \setmainfont[Ligatures={TeX, Common}, Numbers=Lining, ItalicFont=Linux Libertine]{Linux Libertine}
}{}
\IfFontExistsTF{Linux Libertine O}{
  \setmainfont[Ligatures={TeX, Common}, Numbers=Lining, ItalicFont=Linux Libertine O]{Linux Libertine O}
}{}

%% CTeX
% CJK support, used to show CJK characters in the resume
%
% - fontset=none: disable builtin fontset but instead set the CJK font manually
% - heading=false: disable ctex heading
% - punct=kaiming: use kaiming punctuations styles for CJK
% - scheme=plain: use plain scheme, do not override \normalsize font size
% - space=auto: space settings for CJK characters
%
% ref:
% - http://ctan.mirrorcatalogs.com/language/chinese/ctex/ctex.pdf
\usepackage[UTF8, heading=false, punct=kaiming, scheme=plain, space=auto]{ctex}

\IfFontExistsTF{Noto Serif CJK SC}{
  \setCJKmainfont{Noto Serif CJK SC}
}{}
\IfFontExistsTF{Noto Sans CJK SC}{
  \setCJKsansfont{Noto Sans CJK SC}
}{}

%% line spacing
\usepackage{setspace}
\setstretch{1.125}

%% URL styling - use normal text instead of monospace
\urlstyle{same}

% Auto-underline all links
% Defer until \begin{document} so hyperref is already loaded
\AtBeginDocument{%
  \let\oldhref\href
  \renewcommand{\href}[2]{\oldhref{#1}{\underline{#2}}}%
}

%% Patch \httplink and \httpslink to handle full URLs
% The original moderncv macros always prepend http:// or https:// to URLs.
% This causes issues when the URL already has a protocol (e.g., https://example.com)
% resulting in malformed URLs like https://https://example.com.
% This patch checks if the URL already contains "://" and uses it directly if so.
% Note: We use \str_set:Nx (x-expansion) to expand macros like \@homepage first.
\ExplSyntaxOn
\str_new:N \g_yamlresume_url_str

\RenewDocumentCommand{\httpslink}{O{}m}{
  \str_set:Nx \g_yamlresume_url_str {#2}
  \str_if_in:NnTF \g_yamlresume_url_str {://}
    { \url{#2} }
    { \url{https://#2} }
}

\RenewDocumentCommand{\httplink}{O{}m}{
  \str_set:Nx \g_yamlresume_url_str {#2}
  \str_if_in:NnTF \g_yamlresume_url_str {://}
    { \url{#2} }
    { \url{http://#2} }
}
\ExplSyntaxOff

\name{佐藤 美咲}{}
\title{ヒト中心の体験を創造するUXデザイナー}
\phone[mobile]{+81 90 1234 5678}
\email{misaki.sato@example.jp}
\homepage{https://misakidesign.example.jp}

\address{日本、渋谷区、東京都、東京都渋谷区神宮前1-2-3、150-0001}{}{}

\extrainfo{{\small \faLinkedin}\ \href{https://www.linkedin.com/in/misaki-sato}{@misaki-sato} {} {} {} • {} {} {}
{\small \faMedium}\ \href{https://medium.com/@misakisato}{@misakisato}}

\begin{document}

\maketitle

\section{基本情報}

\cvline{}{\begin{itemize}
\item ユーザーリサーチからUIデザイン、プロトタイピング、ユーザビリティテストまでを一貫して担当し、Web・モバイルアプリの体験改善に貢献
\item アクセシビリティとデザインシステムを重視し、エンジニアやプロダクトマネージャーと協働しながら高品質なデザインを提供
\item 定性・定量データを基にした人間中心設計プロセスを実践し、ビジネス目標とユーザー満足度の両立を目指す
\end{itemize}}
\section{学歴}

\cventry{2015年4月～2019年3月}
        {学士、情報デザイン、成績：3.7}
        {多摩美術大学}
        {\url{https://www.tamabi.ac.jp/}}
        {}
        {\begin{itemize}
\item デザイン思考と人間中心設計を学び、ユーザビリティ評価やプロトタイピングの基礎を修得
\item 卒業制作では「視覚障がい者向けナビゲーションアプリ」をデザインし、学内優秀賞を受賞
\end{itemize}
\textbf{コース}：ヒューマンインターフェース、インタラクションデザイン、情報デザイン演習、デザインリサーチ、メディアアート、ユニバーサルデザイン}
\section{職歴}

\cventry{2021年4月～現在}
        {シニアUXデザイナー}
        {株式会社インサイト・デザイン}
        {\url{https://www.insight-design.example.jp}}
        {}
        {\begin{itemize}
\item 大手金融機関のモバイルバンキングアプリのUX改善プロジェクトをリード
\item ユーザーインタビュー、アンケート調査、ヒューリスティック評価を設計・実施し、サービス全体の改善提案に反映
\item デザインシステムの構築と運用を推進し、デザイン品質の一貫性と開発効率を向上
\item クロスファンクショナルチームのファシリテーションを担当し、ステークホルダーとの合意形成を支援
\end{itemize}
\textbf{キーワード}：UXリサーチ、デザインシステム、Figma、ファシリテーション}

\cventry{2019年4月～2021年3月}
        {UXデザイナー}
        {株式会社ウェーブテクノロジー}
        {\url{https://www.wave-tech.example.jp}}
        {}
        {\begin{itemize}
\item ECアプリや社内業務システムのユーザーインターフェース設計に従事
\item ペルソナ設定やカスタマージャーニーマップを作成し、新機能のコンセプト設計を主導
\item プロトタイプを用いたユーザーテストを毎月実施し、改善効果を数値で報告
\end{itemize}
\textbf{キーワード}：ユーザーテスト、プロトタイピング、Adobe XD、情報アーキテクチャ}
\section{言語}

\cvline{日本語}{ネイティブ・母国語レベル \hfill \textbf{キーワード}：母国語}
\cvline{英語}{完全な実務レベル \hfill \textbf{キーワード}：TOEIC 930、ビジネス英会話}
\section{スキル}

\cvline{UXリサーチ}{エキスパート \hfill \textbf{キーワード}：ユーザーインタビュー、ペルソナ設計、カスタマージャーニーマップ、A/Bテスト、アクセシビリティ}
\cvline{UIデザイン}{エキスパート \hfill \textbf{キーワード}：Figma、Adobe XD、Sketch、プロトタイピング、デザインシステム}
\cvline{プロトタイピング・評価}{上級 \hfill \textbf{キーワード}：ユーザビリティテスト、ヒューリスティック評価、リモートユーザーテスト、インサイト分析}
\section{受賞・表彰}

\cventry{2022年10月}
        {公益財団法人日本デザイン振興会}
        {グッドデザイン賞}
        {}
        {}
        {視覚障がい者向けナビゲーションアプリ「まちナビ」のUXデザインが評価され、グッドデザイン賞を受賞。}

\cventry{2019年3月}
        {多摩美術大学}
        {学内優秀賞}
        {}
        {}
        {卒業制作「視覚障がい者向けナビゲーションアプリ」に対して、学内優秀賞を授与。}
\section{資格・免状}

\cventry{2020年11月}
        {一般社団法人人間中心設計推進機構}
        {人間中心設計 基礎認定}
        {\url{https://www.hcdnet.org/}}
        {}
        {}

\cventry{2021年6月}
        {Google}
        {Google UXデザイン プロフェッショナル認定}
        {\url{https://www.coursera.org/professional-certificates/google-ux-design}}
        {}
        {}
\section{出版・著作}

\cventry{2023年3月}
        {ユーザーリサーチを組織に根付かせるための実践ガイド}
        {UXデザイン協会 会誌}
        {\url{https://www.example.jp/articles/ux-research-guide}}
        {}
        {企業におけるユーザーリサーチの定着プロセスを実例に基づき解説。リサーチ結果をプロダクト開発へ活かすためのフレームワークを提案。}
\section{プロジェクト}

\cventry{2022年4月～2022年12月}
        {視覚障がい者向け音声ナビゲーションアプリのUXデザイン}
        {まちナビ}
        {\url{https://www.machinavi.example.jp/}}
        {}
        {\begin{itemize}
\item 視覚障がい者へのインタビューとフィールド調査を実施し、移動中の情報ニーズを明確化
\item 画面読み上げ機能との親和性を考慮した情報設計と音声UIを設計
\item 実機を用いたユーザーテストを繰り返し、操作性と安心感を向上
\end{itemize}
\textbf{キーワード}：音声UI、アクセシビリティ、フィールド調査、ユーザーテスト}
\section{興味・関心}

\cvline{写真撮影}{フィルムカメラ、風景写真}
\cvline{旅行}{温泉めぐり、地域文化}
\section{ボランティア活動}

\cventry{2022年4月～現在}
        {オーガナイザー}
        {UXデザインコミュニティ東京}
        {\url{https://www.ux-tokyo.example.jp}}
        {}
        {\begin{itemize}
\item 若手デザイナー向けの勉強会やポートフォリオレビュー会を企画・運営
\item デザイナー同士の交流を促進し、コミュニティの活性化に貢献
\end{itemize}}

\end{document}